problem:  
Let \(M^5\) be a compact, simply connected Riemannian 5‑manifold with nonnegative sectional curvature that admits an effective, isometric action of the 2‑torus \(T^2\). The orbit space \(X=M/T^2\) is a 3‑dimensional Alexandrov space with nonnegative curvature. It is known (Grove–Searle, Grove–Wilking–Ziller, Oh, and others) that topologically such \(M\) are homeomorphic to one of finitely many known spaces (e.g. \(S^5\), \(S^3\times S^2\), or their connected sums), and that the isotropy graph in \(X\) largely determines the diffeomorphism type.  

In all known examples, different nonnegatively curved metrics produce *different* Alexandrov metrics on their orbit spaces.  

**Concrete Open Rigidity Problem:**  
Does there exist a pair of inequivalent nonnegatively curved metrics \(g_1, g_2\) on the same simply connected 5‑manifold \(M\) with an effective, isometric \(T^2\)-action such that the induced orbit Alexandrov spaces \((X_{g_1},d_{g_1})\) and \((X_{g_2},d_{g_2})\) are *isometric as metric stratified spaces* (same underlying 3‑space with isotropy graphs and link metrics)?  

Equivalently:  
*Can the orbit space metric of a nonnegatively curved 5‑manifold with isometric \(T^2\)‑action fail to determine the equivariant isometry class of the metric?*  

If such a pair does not exist, a natural conjecture follows:  

> **Conjecture (Equivariant Alexandrov Rigidity in Dimension 5):**  
> For simply connected compact \(5\)-manifolds with nonnegative sectional curvature and effective isometric \(T^2\)-actions, the metric on the orbit Alexandrov space, together with isotropy labels, uniquely determines the Riemannian metric up to equivariant isometry.

Why is it a "good" problem:  
This formulation isolates a sharply stated, testable phenomenon: *Is the quotient Alexandrov space a complete metric invariant of the Riemannian geometry under nonnegative curvature and torus symmetry?* It uses well‑understood examples from existing classification results (so the context is rigorous and grounded) but asks a genuinely open, conceptually clean rigidity question—whether distinct global geometries can share identical orbit space metrics.  

Its simplicity (comparing two metrics on the same manifold) belies potentially deep implications: a positive answer would reveal unexpected geometric rigidity linking Alexandrov orbit geometry to smooth curvature data; a negative answer would exhibit hidden moduli of nonnegatively curved metrics invisible to their quotient geometry. Either outcome would have wide consequences for the metric classification program, connecting Alexandrov geometry, torus symmetry, and curvature rigidity.  

The problem works as a focused instance—precision without over‑technicality—illustrating the ideal of “narrow entrance, profound depth.”